% `tex.web` §175: a ligature shows the characters it was *made of*.
%
%     ligature_node: short_display(lig_ptr(p)); {recursive call}
%
% §174's `short_display` is the one-line summary of a box's contents that §663
% prints above an underfull or overfull warning, and §857 prints in a
% paragraph's narration. For a ligature it recurses into the original
% characters, so `office` reads `office` however many ligatures cmr10 made of
% it — `ffi` is one node and prints three letters.
%
% We printed the ligature's own character. Those live where the font keeps
% them, not where the letters do, so `office fluffy waffle` came out as
% `oce uy wae`: every ligature silently *missing* from the line the warning is
% about, which is the line a reader is trying to identify.
%
% In trip.tex the font's ligature for `BB` sits at `A`, so a paragraph of six
% B's read `AAA` — three characters of the wrong letter, and the count wrong
% too.
\catcode`\{=1 \catcode`\}=2
\tracingonline=1
\hbadness=0 \hfuzz=0pt
\showboxbreadth=1 \showboxdepth=0
\hsize=200pt \parindent=0pt
\font\rip=cmr10 \rip

office fluffy waffle\par

% A discretionary's halves go through the same routine, so a hyphenated
% ligature shows its letters on both sides of the break.
\hsize=48pt
office fluffy waffle\par
\end
